My current second bachelor's in Computer Science and Technology has expanded my mechanical background toward algorithms and software. Building a CR5 OpenPI pipeline and modular ROS2 nodes for the HERO RoboMaster team showed me how perception, computation, communication, and control depend on one another. What I still need is deeper training in visual computation and efficient learning, plus a sustained project structure in which performance decisions can be followed through an entire system. CUHK's MSc in Computer Science provides that route.

Where these current courses are available in my intake, CMSC5711 Image Processing and Computer Vision would strengthen the perception side of the pipeline, while CMSC5743 Efficient Computing of Deep Neural Networks would make computational cost part of model and deployment design. With project approval, I would use the CMSC5720 Project I and CMSC5721 Project II sequence to develop one perception-and-policy system across successive stages: establish a visual baseline, improve its computational efficiency, and study how those changes affect serving and robot response. My experience with ROS2 and service interfaces would give the sequence an end-to-end starting point. I want to leave the programme able to build efficient stacks in which sensing, learning, and execution improve together, preparing me for embodied-agent systems R\&D.
